\chapter{Introduction}

\section{Background and Motivation}

Grammatical Error Correction (GEC) is a core task in Natural Language Processing (NLP) concerned with automatically detecting and correcting errors in learner writing. The field has evolved significantly, progressing from rule-based systems in the 1980s and 1990s to statistical machine translation approaches in the early 2000s, and subsequently to neural sequence-to-sequence models in the mid-2010s \cite{10.1145/3474840}. These developments have led to substantial improvements in correction accuracy and scalability, increasing the feasibility of automated writing support in educational contexts.

More recently, large language models (LLMs) have demonstrated strong zero-shot and few-shot performance in GEC tasks \cite{loem2023exploringeffectivenessgpt3grammatical}. Models such as GPT-style systems, Gemini, and LLaMA illustrate the ability of LLMs to perform grammatical correction without task-specific fine-tuning. Unlike earlier approaches, LLMs rely on large-scale generative pre-training and in-context learning, enabling them to produce fluent and contextually coherent revisions \cite{NEURIPS2020_1457c0d6}.

While these capabilities represent a major advancement, they also introduce new challenges. In particular, LLMs tend to prioritise fluency and naturalness, which can result in outputs that go beyond minimal grammatical correction.

\section{Problem Statement}

Rather than performing targeted grammatical edits, LLMs often produce broader rewrites that modify lexical choice, tone, or syntactic structure beyond what is strictly necessary. This phenomenon, referred to as \textit{over-correction}, reflects a bias towards producing fluent and natural-sounding text.

Although such revisions may improve readability, they can introduce unnecessary changes that obscure the original structure of learner writing. For second-language (L2) learners, this behaviour is problematic. Effective feedback should not only correct errors but also preserve the learner's original wording, allowing them to identify and reflect on specific linguistic mistakes.

Excessive rewriting reduces transparency and may limit opportunities for error noticing, which is a key component of language learning. Furthermore, traditional evaluation metrics in GEC, such as ERRANT-based precision, recall, and $F_{0.5}$, focus on agreement with reference corrections and do not fully account for over-correction risk. At the same time, measuring textual change alone is insufficient, because some changes may be useful if they preserve meaning and improve fluency.

These limitations show why GEC methods should aim to correct errors accurately while avoiding unnecessary changes to the learner's original writing.

\section{Aims and Objectives}

The primary aim of this dissertation is to investigate how prompt-based strategies influence potential over-correction in LLM-based grammatical error correction while maintaining effective correction performance.

To achieve this aim, the study pursues the following objectives:

\begin{itemize}
    \item to analyse the nature of over-correction in LLM-generated outputs;
    \item to develop and evaluate multiple prompting strategies, including instruction-constrained, role-based, and few-shot approaches;
    \item to compare prompting strategies using both grammatical accuracy and over-correction risk indicators;
    \item to introduce and apply OCI as an indicator of over-correction risk;
    \item to examine the trade-off between correction performance and over-correction risk;
    \item to investigate whether sentence length influences correction behaviour;
    \item to evaluate the robustness of OCI under alternative weighting schemes.
\end{itemize}

Through these objectives, the dissertation aims to identify prompt designs that reduce over-correction risk while maintaining useful grammatical correction.

\section{Scope of the Study}

This study focuses on evaluating prompt engineering within a controlled experimental setting. A single large language model is used throughout the experiments to ensure that observed differences in output can be attributed to prompt design rather than variation across models.

The evaluation is conducted using the W\&I + LOCNESS dataset, which provides learner-written sentences with annotated corrections. The model is not fine-tuned; instead, all experiments rely on zero-shot and few-shot prompting.

The study emphasises evaluation rather than deployment. Its goal is not to build a production-ready GEC system, but to analyse how different prompting strategies influence correction behaviour, particularly in relation to over-correction.

While the findings provide insights into LLM-based GEC, they are specific to the experimental setup and should not be assumed to generalise to all models or datasets without further validation.

\section{Overview of the Dissertation Structure}

The remainder of this dissertation is organised as follows.\\

\textbf{Chapter 2} presents a critical review of the literature on grammatical error correction (GEC), with particular emphasis on evaluation practices and the limitations of accuracy-based metrics such as ERRANT. It examines the emergence of large language models (LLMs) in GEC, highlighting their strengths in fluency and generalisation alongside their tendency towards over-correction. The chapter further explores over-correction, its pedagogical implications for second-language (L2) learning, and the role of prompt engineering as a mechanism for influencing model behaviour. The chapter concludes by identifying key research gaps, particularly the lack of evaluation frameworks that indicate over-correction risk and the limited empirical study of prompt-based control in learner-oriented settings.\\

\textbf{Chapter 3} translates these gaps into a structured experimental framework. It defines the research questions, dataset selection, and prompting strategies evaluated in this study, including baseline, instruction-constrained, role-based, and few-shot approaches. The chapter outlines the methodology for comparing prompt conditions under a controlled setup, and introduces a dual evaluation framework combining ERRANT-based accuracy metrics with over-correction risk indicators.\\

\textbf{Chapter 4} describes the design and implementation of the experimental pipeline. It details how original learner sentences are processed, how prompts are applied consistently across conditions, and how outputs are stored and evaluated. The chapter explains the implementation of both accuracy-based and change-based evaluation components, including the calculation of OCI. It also outlines the analysis procedures used in the study, including prompt refinement, strategy comparison, sentence-length analysis, robustness testing, trade-off analysis, and qualitative inspection, as well as the steps taken to ensure reproducibility and validity.\\

\textbf{Chapter 5} presents the results and discussion of the experimental study. The analysis is structured around six experiments: prompt refinement, comparison of prompting strategies, sentence-length effects, robustness of the OCI metric, trade-off analysis between correction performance and over-correction, and qualitative analysis of model behaviour. The chapter demonstrates how different prompting strategies produce distinct correction behaviours, and shows that structured prompts can reduce over-correction risk while maintaining strong grammatical performance. It also highlights the relationship between sentence complexity and potential over-correction, and provides qualitative examples to explain how over-correction can manifest in practice.\\

\textbf{Chapter 6} concludes the dissertation by summarising the key findings and contributions of the study. It discusses the implications of prompt engineering for controlling correction behaviour in LLM-based GEC, particularly in educational contexts where minimal and transparent feedback is important. The chapter also outlines the limitations of the study, including dataset and model constraints, and proposes directions for future research, such as adaptive prompting strategies and more fine-grained evaluation of over-correction.\\

% Overall, the dissertation demonstrates that prompt design plays a central role in shaping the balance between grammatical accuracy and over-correction risk, and provides a structured framework for analysing potential over-correction in LLM-based grammatical error correction. 
